% !PRINTED 568
\weberpage{568}

For the prime 17, 3 is a primitive root, and one finds
the index table

\[
\begin{array}{c|cccccccccccccccc}
I & 0 & 1 & 2 & 3 & 4 & 5 & 6 & 7 & 8 & 9 & 10 & 11 & 12 & 13 & 14 & 15\\
\hline
N & 1 & 3 & 9 & 10 & 13 & 5 & 15 & 11 & 16 & 14 & 8 & 7 & 4 & 12 & 2 & 6
\end{array}
.
\]

The two-term periods are, if we set $r = e^{\frac{2\pi i}{17}}$,

\weq{16}{
\begin{aligned}
\eta &= 2\cos\frac{2\pi}{17}, & \eta_1 &= 2\cos\frac{6\pi}{17} \\
\eta_2 &= 2\cos\frac{18\pi}{17} = -2\cos\frac{\pi}{17}, & \eta_3 &= 2\cos\frac{20\pi}{17} = -2\cos\frac{3\pi}{17}, \\
\eta_4 &= 2\cos\frac{26\pi}{17} = 2\cos\frac{8\pi}{17}, & \eta_5 &= 2\cos\frac{10\pi}{17} = -2\cos\frac{7\pi}{17}, \\
\eta_6 &= 2\cos\frac{30\pi}{17} = 2\cos\frac{4\pi}{17}, & \eta_7 &= 2\cos\frac{22\pi}{17} = -2\cos\frac{5\pi}{17}.
\end{aligned}
}

For the computation of the functions $\psi$ one most simply
applies the following table:

\[
\begin{array}{r|ccccccccccccccc}
t & 1 & 2 & 3 & 4 & 5 & 6 & 7 & 8 & 9 & 10 & 11 & 12 & 13 & 14 & 15\\
\hline
\operatorname{ind} t & 0 & 14 & 1 & 12 & 5 & 15 & 11 & 10 & 2 & 3 & 7 & 13 & 4 & 9 & 6 \\
\hline
\operatorname{ind}(t+1) & 14 & 1 & 12 & 5 & 15 & 11 & 10 & 2 & 3 & 7 & 13 & 4 & 9 & 6 & 8
\end{array}
,
\]

from which, if $\alpha$ denotes an 8th root of unity, one finds by
formula (8):

\weq{17}{
\begin{aligned}
\psi_1(\alpha) &= \psi_6(\alpha) = 2\alpha + 2\alpha^2 + 3\alpha^4 + 4\alpha^5 + 2\alpha^6 + 2\alpha^7 \\
\psi_2(\alpha) &= \psi_5(\alpha) = 2 + 3\alpha + \alpha^3 + \alpha^4 + 3\alpha^5 + 4\alpha^6 + \alpha^7 \\
\psi_3(\alpha) &= \psi_4(\alpha) = 3 + 3\alpha + 2\alpha^2 + 3\alpha^3 + \alpha^5 + 2\alpha^6 + \alpha^7.
\end{aligned}
}

If one takes

\[
\alpha = \frac{1+i}{\sqrt{2}}, \qquad \alpha^2 = i, \qquad \alpha^4 = -1,
\]

there results

\weq{18}{
\begin{aligned}
\psi_1(\alpha) &= \psi_6(\alpha) = -3 - i\sqrt{8} \\
\psi_2(\alpha) &= \psi_5(\alpha) = 1 - 4i \\
\psi_3(\alpha) &= \psi_4(\alpha) = 3 + i\sqrt{8} \\
\psi_1(i) &= -1 + 4i,
\end{aligned}
}

hence by (15)

\weq{19}{
(\alpha, \eta)^8 = 17\,(3 + i\sqrt{8})^4\,(1 - 4i)^2.
}

But from (13) and (14) one obtains
%CONT>>
